\chapter{فهم بوزون هيغز}

\lectureoverview{تُفهم آلية هيغز على أفضل وجه وفق الترتيب الذي تصبح فيه مكوّناتها ضرورية. نبدأ بالحركة المكمّاة والحقول، ونحوّل وعاءً عاديًا إلى كمون قبعة مكسيكية، ونستخدم المكثّف الناتج لفهم سبب إمكان امتلاك الفرميونات والبوزونات الضعيفة كتلة. وعندئذ فقط نحدّد بوزون هيغز الفيزيائي بوصفه الاهتزاز الشعاعي لذلك المكثّف ونتتبّع عمليات الإنتاج والاضمحلال التي كشفت عنه في 2012.}

\section{الآلية الكامنة وراء العناوين الرئيسية}

كان اكتشاف جسيم جديد قرب $125\,\mathrm{GeV}$ مثيرًا بحق، لكن الإثارة ليست فهمًا بعد. والسؤال المفيد ليس ما إذا كان هيغز مهمًا، ولا ما إذا كان يمكن وصله بلاغيًا بأصل الكون. بل السؤال المفيد أضيق بكثير: ما المكوّنات العملية لهذه الظاهرة، ولماذا احتاج إليه النموذج المعياري؟

ويستحق الاسم أيضًا تحفّظًا موجزًا. فقد كان Peter Higgs واحدًا من عدة أشخاص حدّدوا هذه الفيزياء، كما توصّل Fran\c{c}ois Englert وRobert Brout إلى الفكرة المركزية على نحو مستقل. ولذلك فإن عبارة \emph{آلية Brout--Englert--Higgs} تسجّل التاريخ على نحو أوفى، مع أن الاسم الأقصر سيُستخدم أدناه.

يأتي الشرح في وحدات. نحتاج أولًا إلى قطعتين أوليتين من ميكانيكا الكم وإلى فكرة واحدة من نظرية الحقول الكلاسيكية. ثم نحتاج إلى فراغ تحتوي حالته ذات الطاقة الدنيا على حقل غير صفري. وبعد ذلك يمكننا أن نجرد جسيمات النموذج المعياري ونرى لماذا تترك نسخة عارية من النظرية الجسيمات المعنية بلا كتلة. أما بوزون هيغز نفسه فلا يظهر إلا قرب النهاية. فهو ليس الآلية كلها؛ بل هو اهتزاز معيّن واحد يتبقى بعد أن تصبح الآلية قائمة.

\section{حقيقتان كمّيتان وحقل}

الحقيقة الكمّية الأولى هي التكميم. ويُعدّ الزخم الزاوي نموذجًا أوليًا واضحًا لأنه يأتي في وحدات منفصلة،
\begin{equation}
L=n\hbar,
\qquad
n=0,\pm1,\pm2,\ldots .
\end{equation}
وتوجد أيضًا عزوم زاوية نصف-صحيحة، لكن لا شيء في الحجة المباشرة يعتمد عليها. والمهم هو أن حركة كلاسيكية تتغير على نحو متصل تصبح سُلّمًا من الحالات الكمّية.

الحقيقة الثانية هي مبدأ عدم اليقين. وسنؤجلها إلى أن تظهر دائرة متصلة من الفراغات الممكنة. وعندئذ سيمنعنا من معاملة الفراغ كما لو كان كرة كلاسيكية صغيرة ساكنة عند نقطة واحدة محددة تحديدًا حادًا.

بين هاتين الحقيقتين الكمّيتين يقع المفهوم الكلاسيكي للحقل. فالحقل يعيّن كمية لكل نقطة من المكان والزمان. والحقول الكهربائية والمغناطيسية والثقالية أمثلة مألوفة. وقد تبدو منطقة ما فارغة مع أنها لا تزال تحتوي على حقل. تخيّل لوحين مشحونين متباعدين إلى حد لا يمكن معه رؤية أيٍّ منهما. ومع ذلك فإن الحقل الكهربائي شبه المنتظم بينهما حقيقي محليًا. وهكذا فإن ``الفضاء الفارغ'' لا يلزم أن يعني ``أن كل حقل يساوي الصفر.'' بل يعني أن الحقول تشغل حالتها ذات أدنى طاقة متاحة.

افترض أن طاقة مجال ما تكون أصغر ما يمكن عند قيمة واحدة، وتزداد عندما يُزاح المجال عنها. وإذا أُحدث اضطراب في المجال، فإنه يتذبذب حول القيمة الصغرى. وتُكمِّم ميكانيكا الكم تلك التذبذبات، وتكون الكمّات المفردة جسيمات. هذا هو القاموس اللازم في نظرية الحقول طوال المحاضرة:
\begin{equation}
\text{particle}
\quad\longleftrightarrow\quad
\text{quantum of a field vibration}.
\end{equation}

\section{وعاء في فضاء المجال}

بالنسبة إلى مركبة مجال حقيقية واحدة $\phi$، فإن أبسط شكل لتضاريس الطاقة هو وعاء،
\begin{equation}
V(\phi)=\frac{1}{2}m^{2}\phi^{2}.
\end{equation}
ويقع قاع الوعاء عند $\phi=0$. وللاهتزازات الصغيرة حول القاع تردد محدد، وتتصرف كمّاتها بوصفها جسيمات كتلتها $m$.

الآن استبدل المجال المفرد بمركبتين، $(\phi_1,\phi_2)$. ولا يمكن أن تعتمد الطاقة إلا على نصف القطر في هذا المستوى الداخلي،
\begin{equation}
r^{2}=\phi_1^{2}+\phi_2^{2},
\qquad
V=V(r).
\end{equation}
هذا وعاء فوق \emph{فضاء المجال}، لا انخفاضًا في الفضاء العادي. إن تدوير الزوج $(\phi_1,\phi_2)$ يترك الطاقة من دون تغيير. والقبعة المستخدمة في المحاضرة تجعل التناظر نفسه مرئيًّا: أدرها حول محورها فستظل تبدو كما هي.

\begin{figure}[tbp]
\centering
\includegraphics[width=0.95\textwidth]{lecture_144_figure_02.png}
\caption{الوعاء المتماثل العادي والقبعة المادية المستخدمة لجعل التناظر الدوراني في فضاء المجال ملموسًا.}
\lectureframe{00:10:02}
\end{figure}

تصبح المماثلة أكثر من مجرد هندسة عندما يتحرك المجال. أزِحه عن المركز واجعله يدور حول الوعاء. وتحمل هذه الحركة الدورانية زخمًا زاويًّا في فضاء المجال. وليس هذا هو الزخم الزاوي لجسم يدور في المختبر؛ بل هو زخم زاوي داخلي. وعند تكميمه، فإنه يوفّر نموذجًا لشحنة محفوظة،
\begin{equation}
Q=nq_0,
\qquad
n\in\mathbb{Z}.
\end{equation}
والمعادلة تخطيطية. والفكرة الفيزيائية المركزية هي أنه يمكن تصوير إثارة مجال مشحونة على أنها دوران في مستوى داخلي. وتقابل الشحنة الموجبة والسالبة اتجاهين متعاكسين للدوران، ويقابل انعدام الدوران شحنة صفرية.

\section{قلب القبعة}

للوعاء العادي نقطة دنيا واحدة فقط. وتبدأ فيزياء هيغز عندما تُقلَب القبعة. فالمركز الآن غير مستقر، في حين أن أخدودًا دائريًّا كاملًا له الطاقة الدنيا نفسها. وأحد التحققات الرياضية المعيارية لذلك هو
\begin{equation}
V(\phi_1,\phi_2)
=
\lambda
\left(
  \phi_1^{2}+\phi_2^{2}-v^{2}
\right)^{2},
\qquad \lambda>0,
\end{equation}
وتحقق نقاطه الدنيا
\begin{equation}
\phi_1^{2}+\phi_2^{2}=v^{2}.
\end{equation}
\editorialnote{إن الجهد الرباعي الصريح هو الإتمام التحليلي المعياري لهندسة القبعة المكسيكية المرسومة في المحاضرة؛ وليس منسوخًا من صيغة على اللوح.}

\begin{figure}[tbp]
\centering
\includegraphics[width=0.95\textwidth]{lecture_144_figure_03.png}
\caption{القمة المركزية والأخدود الدائري لجهد القبعة المقلوبة. يحل الرسم الكامل على اللوح محل اقتطاع سابق كان قد حذف معظم الهندسة.}
\lectureframe{00:13:10}
\end{figure}

لم يعد الخلاء عند مجال صفري. بل لا بد أن يقع في مكان ما على الأخدود،
\begin{equation}
(\phi_1,\phi_2)
=
v(\cos\theta,\sin\theta).
\end{equation}
وتحتفظ المعادلات بالتناظر الدوراني، لكن أي خلاء بعينه يختار زاوية واحدة $\theta$. وهذا هو الكسر التلقائي للتناظر: فالقانون متناظر، في حين أن حالة الطاقة الدنيا المختارة ليست كذلك.

يوجد الآن نوعان مختلفان نوعيًا من الحركة. فالتحرك شعاعيًا إلى الداخل أو إلى الخارج يصعد في الجهد ويتطلب طاقة. أما التحرك حول الحوض فلا يغيّر إلا $\theta$ ولا يكلّف، على المستوى الكلاسيكي، أي طاقة وضع. ولذلك يمكن أن تقترن كمية كبيرة من الحركة الزاوية الشبيهة بالشحنة بقدر ضئيل جدًا من طاقة الوضع الإضافية.

\begin{figure}[tbp]
\centering
\includegraphics[width=0.95\textwidth]{lecture_144_figure_04.png}
\caption{تمثّل نقطة معلَّمة في الحوض فراغًا معيّنًا مختارًا من الدائرة المتصلة لنقاط الحد الأدنى المتكافئة جميعًا.}
\lectureframe{00:16:14}
\end{figure}

\section{يحوّل عدم اليقين الفراغ إلى مُكثَّف}

كلاسيكيًا، قد يضع المرء الحقل عند نقطة على الحافة، ويضبط سرعته الزاوية على الصفر، ويعلن أن المسألة قد انتهت. لكن ميكانيكا الكم تمنع مثل هذا الحسم الحاد. فالموضع والزخم يحققان
\begin{equation}
\Delta x\,\Delta p\gtrsim\frac{\hbar}{2}.
\end{equation}
والعبارة المناظرة على الحوض الدائري هي أن حصر زاوية الحقل يجعل زخمه الزاوي الداخلي المرافق غير يقيني. وبلغة المخطط في المحاضرة،
\begin{equation}
\Delta\theta\,\Delta L_{\mathrm{int}}
\gtrsim
\frac{\hbar}{2}.
\end{equation}
\editorialnote{لمؤثرات الزاوية دقائق تقنية، لذا فهذه العلاقة تعبير محلي شبه كلاسيكي عن حجة عدم اليقين في المحاضرة، وليست هوية مؤثرات مضبوطة على نحو تام عالميًا.}

لذلك فإن الفراغ المتركز قرب نقطة واحدة من الحوض ليس حالةً ذات شحنة داخلية محددة. فهو يحتوي تراكبًا واسعًا من الحالات الذاتية للشحنة. وتُحكِم المحاضرة هذا المعنى باختبار إجرائي: فإضافة وحدة واحدة من الشحنة، أو إزالة وحدة واحدة، لا تغيّر الحالة إلا بالكاد. ويمكن تصوير ذلك تخطيطيًا كما يلي
\begin{align}
\ket{\Omega}
&\sim
\sum_n c_n\ket{n},
\\
c_{n+1}
&\simeq
c_n
\quad\text{within the occupied range}.
\end{align}
وعندئذ فإن إزاحة $n$ بمقدار واحد تعيد في معظمها تسمية حدود موجودة أصلًا. وهذا هو المعنى الذي يكون فيه الفراغ \emph{مُكثَّفًا}: فهو يحتوي عددًا غير محدد من الكمّات ذات الصلة، بحيث إن كمًّا واحدًا أكثر أو كمًّا واحدًا أقل لا ينتج حالة عيانية مختلفة على نحو قابل للرصد.

الفراغ الكهروضعيف الحقيقي ليس مُكثَّفًا من الشحنة الكهربائية العادية. ولو كان كذلك، لكان السلوك الكهرومغناطيسي للفضاء الخالي مختلفًا اختلافًا جذريًا. غير أن الموصل الفائق يقدم مثالًا ملموسًا على مُكثَّف مشحون. ولذلك فهو يشكّل جسرًا فيزيائيًا مفيدًا بين صورة القبعة المكسيكية المجردة والقول بأن الفراغ يمكن أن يحمل بنية كمومية متماسكة.

بهذا تكتمل الوحدة الأولى. وقد تعلمنا أن حالة أدنى طاقة لا يلزم أن تقع عند انعدام الحقل، وأن عدم اليقين الكمومي يمكن أن يجعل تلك الحالة غير حساسة لإضافة الكمّات المفردة أو إزالتها. ونسأل الآن أيُّ الجسيمات تحتاج إلى مثل هذه الخلفية.

\section{النموذج المعياري العاري}

المعيار العلوي الطبيعي لكتل الجسيمات الأولية هو كتلة بلانك. فإذا أوصلت طاقةً متموضعة إلى حد عالٍ بما يكفي، لم يعد بالإمكان إهمال الجاذبية؛ وفي النهاية يغدو من الأنسب عدُّ الجسم ثقبًا أسود بدلًا من جسيم أولي عادي. عدديًا،
\begin{equation}
m_{\mathrm P}
=
\sqrt{\frac{\hbar c}{G}}
\approx
2.2\times10^{-8}\,\mathrm{kg},
\end{equation}
أي تقريبًا كتلة ذرة غبار صغيرة. والجسيمات الأولية المعروفة أخف بدرجة مذهلة. وفي تقدير السُّلَّم المستخدم في المحاضرة، لا تبلغ كتلها إلا نحو $10^{-17}$ من كتلة بلانك.

هذه الملاحظة لا تثبت عدم وجود جسيمات أولية أثقل. فالمسرِّعات تكشف على نحو تفضيلي الجسيمات الخفيفة بما يكفي لإنتاجها بالطاقة المتاحة. ولذلك فإن الطيف المعروف مُنتقى جزئياً بمدى الآلات المستخدمة للبحث عنه. وقد توجد جسيمات أثقل كثيرة تتجاوز إمكان الوصول الحالي.

تنقسم الجسيمات المرصودة إلى فرميونات وبوزونات. وتشمل الفرميونات الإلكترون والميون والتاو ونيوترينواتها ونكهات الكواركات الست. وتشمل البوزونات الفوتون والغلوونات والبوزونات الضعيفة $W$ و$Z$. وتربط الغلوونات الكواركات داخل البروتونات والنيوترونات، ولذلك فهي عنصر محوري في بنية النوى العادية.

\begin{figure}[tbp]
\centering
\includegraphics[width=0.95\textwidth]{lecture_144_figure_05.png}
\caption{الجرد الأول للنموذج المعياري في المحاضرة، مع رسم المدى الهائل للكتل فوق عمودي الفرميونات والبوزونات.}
\lectureframe{00:24:22}
\end{figure}

ليست النظرية مجرد قائمة؛ بل هي مجموعة من التفاعلات المسموح بها. يمكن للإلكترون أن يصدر فوتوناً،
\begin{equation}
e\longrightarrow e+\gamma.
\end{equation}
ويحمل الكوارك كلاً من الشحنة الكهربائية وشحنة اللون، لذا يمكنه أن يصدر إما فوتوناً أو غلووناً،
\begin{equation}
q\longrightarrow q+\gamma,
\qquad
q\longrightarrow q+g.
\end{equation}
ولا يستطيع الإلكترون أن يصدر غلووناً، كما أن النيوترينو لا يصدر لا فوتوناً ولا غلووناً. ويمكن للفرميونات ذات الأعداد الكمومية الضعيفة المناسبة أن تتفاعل عبر $Z$،
\begin{equation}
f\longrightarrow f+Z.
\end{equation}
كما أن بوزونات $W$ وتفاعلات التيار المشحون جزء من النموذج المعياري أيضاً، لكن المحاضرة تُغفلها لكي يتولى بوزون ضعيف متعادل واحد حمل الحجة.

\begin{figure}[tbp]
\centering
\includegraphics[width=0.95\textwidth]{lecture_144_figure_06.png}
\caption{رؤوس إصدار الفوتون والغلوون و$Z$ بجوار كلمة ``عديم الكتلة''. تلخّص هذه العمليات النظرية العارية قبل إدخال المتكاثف.}
\lectureframe{00:29:20}
\end{figure}

الفوتون عديم الكتلة. وفي نظرية معيارية غير مكسورة من النوع العام نفسه، تكون الغلوونات والبوزونات المعيارية الضعيفة أيضاً عديمة الكتلة. وكذلك لا تستطيع النظرية الكهروضعيفة غير المعدلة أن تُدخل ببساطة كتل الفرميونات يدوياً مع الحفاظ على تناظرها المعياري. ثمة عنصر مفقود. وقبل تقديمه، نحتاج إلى أن نرى في سياق أبسط كيف يمكن لحقل خلفي أن يغيّر كتلة جسيم من دون أن يعمل كاحتكاك.

\section{كيف يمكن لحقلٍ خلفي أن يغيّر الكتلة}

تأمل جزيء ماء بوصفه ثنائياً قطبياً كهربائياً: أحد طرفيه موجب قليلاً والآخر سالب قليلاً. وفي الفضاء الخالي، يكون اتجاهان متعاكسان متكافئين من حيث الطاقة. وتمنح المحاضرة هذين الاتجاهين، على سبيل المداعبة، اسمين مختلفين، ``ماء'' و``سكوتش''، لمجرد أن يمكن تتبعهما بوصفهما حالتين داخليتين متميزتين.

ضع الجزيء بين لوحي مكثف. إذا كان عزم ثنائي القطب له هو $\bm d$ وكان المجال الكهربائي هو $\bm E$، فإن طاقة التفاعل هي
\begin{equation}
U=-\bm d\mathbin{\cdot}\bm E.
\end{equation}
وبالنسبة إلى الاتجاهين المحاذيين للمجال يصبح هذا
\begin{equation}
U_{\parallel}=-dE,
\qquad
U_{\mathrm{anti}}=+dE.
\end{equation}
وتحوّل علاقة أينشتاين $E=mc^2$ انقسام الطاقة إلى انقسام في الكتلة،
\begin{equation}
m_{\mathrm{anti}}-m_{\parallel}
=
\frac{2dE}{c^{2}}.
\end{equation}
لم يحتكّ بالجزيء شيء. فإحدى الحالتين الداخليتين تملك ببساطة طاقة أكبر، ومن ثم كتلة أكبر، من الأخرى ما دام الحقل الخلفي موجوداً.

\begin{figure}[tbp]
\centering
\includegraphics[width=0.95\textwidth]{lecture_144_figure_07.png}
\caption{اتجاهان جزيئيان متعاكسان في مجال المكثف. يفصل الحقل بين طاقتيهما من دون أن يُدخل قوة ممانعة.}
\lectureframe{00:34:35}
\end{figure}

وهذا يبدد تشبيهين مضللين. فحقل هيغز ليس دبساً يتعين على الجسيمات أن تشق طريقها عبره، وليس الجسيم ذو الكتلة كاسحةَ ثلج تُبطئها مادة تتراكم أمامها. وفي حقل منتظم مكانياً قد لا توجد أي قوة محصلة أصلاً. فالأثر هو إزاحة في الطاقة الداخلية، وليس احتكاكاً ولا فقداناً للزخم.

ويمكن إعادة سرد المثال نفسه بلغة كمومية. فالمجال الكهربائي الكلاسيكي حالة متماسكة، يمكن وصفها لهذا الغرض بأنها متكاثف من الفوتونات. وقد يصدر الجزيء فوتوناً إلى تلك الخلفية أو يمتص فوتوناً منها. وبما أن إضافة كم واحد أو إزالته لا تغيّر إلا تغييراً طفيفاً حالةً متماسكة عالية الإشغال، فيمكن لخط الفوتون أن ينتهي في المتكاثف. وتشير العلامة المتقاطعة المرسومة على اللوح إلى هذا الاختفاء في الخلفية، لا إلى إخفاق في حفظ الطاقة.

\begin{figure}[tbp]
\centering
\includegraphics[width=0.95\textwidth]{lecture_144_figure_08.png}
\caption{أعيد رسم مجال المكثف بوصفه متكاثفاً فوتونياً. يُمتص فوتون يصدره ثنائي القطب في خلفية لا تتأثر تقريباً بكم إضافي واحد.}
\lectureframe{00:38:45}
\end{figure}

هذا هو القالب الذي نحتاج إليه. فإذا أمكن لمتكاثف أن يمتص فرق الأعداد الكمومية بين حالتين داخليتين، أمكنه أن يخلط بينهما. وقد يظهر التذبذب الناجم عن ذلك على هيئة كتلة.

\section{لماذا لا يفسّر هيغز كل الكتلة}

قبل تطبيق ذلك القالب، يجب إزالة تصور خاطئ واحد. فمن غير الصحيح أن كل جسم ذي كتلة يستمد كل كتلته من حقل هيغز.

ضع إشعاعاً في صندوق عاكس تماماً. الفوتونات المفردة عديمة الكتلة، لكن الصندوق المغلق له إطار سكون ويخزن طاقة. وتزداد كتلته بمقدار
\begin{equation}
\Delta M=\frac{E_{\mathrm{radiation}}}{c^{2}}.
\end{equation}
وبذلك يمكن أن تنشأ الكتلة من طاقة حركية وطاقة مجال محصورتين، حتى عندما لا يكون للمكوِّنات كتلة سكون.

والبروتون هو المثال الواقعي المهم. فهو يحتوي على كواركات وكواركات مضادة وغلوونات في حالة تآثر قوي. وتأتي معظم كتلة البروتون من طاقتها الحركية وطاقة ربطها، لا من مجرد جمع الكتل التي يولدها هيغز لبضعة كواركات مكوِّنة. وفي التقدير الذي تؤكد عليه المحاضرة، فإن جعل الكواركات الخفيفة عديمة الكتلة لن يغيّر كتلة البروتون إلا على مستوى يقارب الواحد في المئة. وتُبرز الثقوب السوداء هذه النقطة المنطقية بصورة أشد وضوحاً: فقد يكون الثقب الأسود ذا كتلة بسبب الطاقة الكلية التي كوّنته، لا لأن كل مساهمة حُوِّلت أولاً بواسطة حقل هيغز إلى كتلة سكون لجسيم أولي.

لآلية هيغز وظيفة أدق. فهي تتيح كتلًا للفرميونات الأولية وللبوزونات المعيارية الضعيفة مع الحفاظ على اتساق النظرية المعيارية الكامنة.

\section{الكيرالية: الكتلة بوصفها مزجاً بين اليسار--اليمين}

لنركز أولاً على الإلكترون. فاللف المغزلي عزم زاوي ذاتي. وبالنسبة إلى جسيم متحرك، يمكن مقارنة اتجاه لفه المغزلي بزخمه. وفي اللغة الهندسية للمحاضرة، يُميَّز الإلكترون اليميني والإلكترون اليساري بالعلاقة بين هذين الاتجاهين. وبالنسبة إلى جسيم عديم الكتلة، تتوافق اللولبية والكيرالية على النحو البسيط اللازم للحجة؛ أما بالنسبة إلى جسيم ذي كتلة، فيجب التمييز بين المفهومين بدقة.

تخيل أن تتجاوز مراراً إلكتروناً ذا كتلة. ففي مرجع قصوري آخر ينعكس زخمه بينما لا ينعكس لفه المغزلي، ومن ثم يمكن أن تنعكس اليدوية الظاهرية المرتبطة باتجاه الحركة. ولا يستطيع أي راصد تجاوز جسيم عديم الكتلة تماماً، لأنه يتحرك دائماً بسرعة الضوء. وهذا يحفّز العبارة الاسترشادية في المحاضرة: فالإلكترون عديم الكتلة يبقى في حالة كيرالية واحدة، في حين أن الإلكترون ذا الكتلة يمزج باستمرار بين اليسار واليمين.

والصياغة الدقيقة في نظرية الحقول هي أن حد كتلة ديراك يقرن الحقلين الكيراليين،
\begin{equation}
\mathcal L_m
=
-m_e\bar e e
=
-m_e
\left(
  \bar e_L e_R
  +
  \bar e_R e_L
\right).
\end{equation}
\editorialnote{حد كتلة ديراك هو الاستكمال الرياضي القياسي لصورة المحاضرة عن ``معدل الانقلاب''. ويتناسب مزج الكيرالية طردياً مع الكتلة؛ أما المعدل الزمني الحرفي فيعتمد على الحالة وعلى تمدد الزمن النسبي.}

إذا كان $m_e=0$ فإن الحقلين الأيسر والأيمن ينتشران على نحو مستقل. وإذا كان $m_e\ne0$ فإن المعادلات تقرن بينهما. وبهذا المعنى الدقيق، تقيس الكتلة شدة المزج بين اليسار--اليمين.

\section{قانون الحفظ الذي يمنع كتلة عارية}

لماذا لا نكتب ببساطة حد المزج منذ البداية؟ يعامل التفاعل الضعيف $e_L$ و$e_R$ على نحو مختلف. ويطلق سسكند على هذا التمييز ذي الصلة اسماً بسيطاً عمداً: \emph{زيلتش}. في التطبيع المستخدم في المحاضرة،
\begin{equation}
z(e_L)=1,
\qquad
z(e_R)=0.
\end{equation}
للكيراليتين الشحنة الكهربائية نفسها، لكن لهما أعداداً كمومية ضعيفة مختلفة. ولذا فإن انتقالاً مباشراً
\begin{equation}
e_L\longrightarrow e_R
\end{equation}
لن يحفظ الكمية التي يحملها التفاعل الضعيف. وفي النظرية غير المكسورة، يكون حد الكتلة العارية ممنوعاً.

\editorialnote{``زيلتش'' اختزال تربوي. في النموذج المعياري، تنتمي الفرميونات اليسارية واليمينية إلى تمثيلات مختلفة لـ $SU(2)_L\times U(1)_Y$، ولذلك فإن حد كتلة ديراك المباشر لا يكون ثابتاً معيارياً. وتُشفِّر تعيينات المحاضرة $1$ و$0$ هذا الاختلاف من دون إعادة إنتاج جدول الشحنات الكهروضعيفة الكامل.}

يعرض اللوح هذا العائق بوصفه سلسلة. فلا يستطيع إلكترون يميني أن يصبح يسارياً ما لم يحمل شيء آخر فرق الزيلتش بعيداً. وينطبق العائق نفسه على الكواركات والميونات والتاوات، مع شدات تعتمد على النوع حين يُزوَّد الحقل المفقود.

\begin{figure}[tbp]
\centering
\includegraphics[width=0.95\textwidth]{lecture_144_figure_09.png}
\caption{التسلسل الكيرالي على اللوح: إدراج كمٍّ من المتكاثف عند علامة التقاطع الخضراء يتيح لوسم الزيلتش أن يتغير بين انتشار يميني وآخر يساري.}
\lectureframe{00:50:15}
\end{figure}

\section{متكاثف ``زيغز'' وكتلة الفرميونات}

لنُدخل حقلاً تحمل كمّاته بالضبط عدد الكم الضعيف المفقود. ولمنع حدوث التباس مبكر مع جسيم هيغز الفيزيائي، تسمّي المحاضرة الكمَّ المعني فيه \emph{Ziggs}. افترض أن لحقل Ziggs جهدَ القبعة المكسيكية وأن فراغه متكاثف. عندئذٍ يستطيع إلكترون أعسر أن يبعث كمًّا من Ziggs ويصبح أيمنًا،
\begin{equation}
e_L\longrightarrow e_R+\text{Ziggs},
\end{equation}
أو يستطيع إلكترون أيمن أن يمتص واحدًا ويصبح أعسرًا،
\begin{equation}
e_R+\text{Ziggs}\longrightarrow e_L.
\end{equation}
ويمتص المتكاثف فرق عدد الكم أو يزوّده من غير أن يغيّر حالته العيانية. وما كان ممنوعًا في الفراغ الخالي غير المنكسر يصبح مسموحًا به في الفراغ ذي التناظر المكسور.

يؤدي الإصدار والامتصاص المتكرران إلى مزج $e_L$ و$e_R$. ولذلك ينتشر الإلكترون بوصفه جسيمًا ذا كتلة. وتنطبق الآلية نفسها على اللبتونات المشحونة والكواركات. ولكل نوع اقترانه الخاص بالمتكاثف، لذا يمكن أن تختلف شدة التقليب، ومن ثم الكتلة، من نوع إلى آخر.

في الترميز الحديث، يكون التفاعل غير المتغيّر تحت التحويلات العيارية اقترانَ يوكاوا. وبعد أن يكتسب حقل هيغز قيمة التوقع في الفراغ $v$، فإنه يولِّد
\begin{equation}
\mathcal L_Y
\supset
-\frac{y_f}{\sqrt{2}}
(v+h)\bar f f,
\end{equation}
ومن ثم
\begin{equation}
m_f=\frac{y_fv}{\sqrt{2}},
\qquad
g_{hff}=\frac{y_f}{\sqrt{2}}
=
\frac{m_f}{v}.
\end{equation}
\editorialnote{توفّر هذه الصيغ الترميز الكهروضعيف المعياري الكامن وراء حجة المحاضرة بشأن اقتران المتكاثف. كما تُظهر صراحةً لماذا يتحكم الاقتران نفسه في كلٍّ من كتلة الفرميون واضمحلال هيغز إلى ذلك الفرميون.}

وتصف المحاضرة هذا بأنه كسر تلقائي لتناظر كيرالي. وهذا يلتقط قصة اليسار--اليمين على مستوى المثال. وعلى نحو أدق، تكسر قيمة توقع هيغز تناظر العيار الكهروضعيف إلى الزمرة الفرعية الكهرومغناطيسية، ثم تصبح تفاعلات يوكاوا حدودَ كتلة للفرميونات.

\section{كيف يكتسب البوزون الضعيف كتلة}

ويغيّر المتكاثف أيضًا بوزونات العيار الضعيفة. فإثارة Ziggs تستطيع أن تبعث بوزونًا من نوع $Z$، وبالعكس يمكن للبوزون $Z$ أن يتحول إلى الإثارة المناسبة المدعومة بالمتكاثف ثم يعود من جديد. وهكذا ينتشر $Z$ عبر فراغ يمتزج فيه بالقطاع الزاوي من حقل هيغز. والنتيجة بوزون متجه ذو كتلة، له استقطاب طولي لا يملكه بوزون متجه عديم الكتلة.

\begin{figure}[tbp]
\centering
\includegraphics[width=0.95\textwidth]{lecture_144_figure_10.png}
\caption{تسلسل اللوح المكتمل لآلية البوزون الضعيف: تبقى السلسلة الكيرالية في الأعلى، بينما يتبادل $Z$ عدد الكم مع القطاع المدعوم بالمتكاثف في الأسفل.}
\lectureframe{00:53:40}
\end{figure}

هذه هي ظاهرة Brout--Englert--Higgs في صورتها الخاصة ببوزونات العيار. وفي النموذج المعياري الكامل، تكتسب بوزونات $W$ والبوزون $Z$ كتلة، بينما يبقى أحد توافيق حقول العيار المتعادلة الأصلية هو الفوتون العديم الكتلة. ولو كان الفراغ، بدلًا من ذلك، قد كسر تناظر العيار الكهرومغناطيسي العادي، لأصبح الفوتون ذا كتلة، كما يحدث له فعليًا داخل موصل فائق. ولكان للفضاء الخالي عندئذٍ سلوك كهرومغناطيسي بعيد المدى مختلف تمامًا.

لقد قدّمت البوزونات الضعيفة بالفعل دليلًا قويًا على هذه الآلية قبل وقت طويل من عزل جسيم هيغز الفيزيائي. وقد أظهرت كتلها وتفاعلاتها أن التناظر الكهروضعيف قد انكسر بالطريقة المطلوبة. ولهذا تستطيع المحاضرة أن تقول، بمعنًى إجرائي فضفاض، إن قطاع ``Ziggs'' كان معروفًا بالفعل: فالشيء التجريبي المفقود كان الاهتزاز الشعاعي المستقل.

\editorialnote{إن الإشارة الارتجالية في المحاضرة إلى اكتشاف في SLAC نحو 1980 تختزل عدة أحداث. فقد رُصدت التيارات الضعيفة المتعادلة في 1973، وأُبلغ عن تشتت الإلكترونات المنتهك لتناظر التكافؤ في SLAC في أواخر 1970يات، واكتُشف بوزونا $W$ و$Z$ مباشرةً في CERN في 1983. والمغزى الفيزيائي هو التأكيد التجريبي السابق لانكسار التناظر الكهروضعيف.}

\section{ما هو بوزون هيغز فعليًا}

لنعد إلى القبعة المكسيكية. فالحركة حول الحضيض تغيّر الزاوية لا نصف القطر. وهذا القطاع الزاوي يزوّد الأنماط التي تُمتص في البوزونات الضعيفة الضخمة. أما بوزون هيغز الفيزيائي فهو الحركة الأخرى: اهتزاز إلى الداخل وإلى الخارج، صعودًا على جدران الجهد.

اكتب المجال تخطيطيًا كما يلي
\begin{equation}
\Phi(x)
=
\frac{v+h(x)}{\sqrt{2}}
\exp\!\left(
  \frac{i\theta(x)}{v}
\right).
\end{equation}
يصف المتغير $\theta$ الحركة على طول الحضيض. ويغيّر المتغير $h$ نصف القطر. وفي نظرية معيارية تصبح درجات الحرية الزاوية المناسبة المركبات الطولية للبوزونات الضعيفة الضخمة، بينما يبقى $h$ جسيمًا قياسيًا منتشرًا: بوزون هيغز.

\begin{figure}[tbp]
\centering
\includegraphics[width=0.95\textwidth]{lecture_144_figure_11.png}
\caption{يظهر الجهدان من جديد على اللوح، مع تحديد إزاحة شعاعية في حضيض القبعة المكسيكية. وهذا الاهتزاز إلى الداخل وإلى الخارج هو نمط هيغز الفيزيائي.}
\lectureframe{00:56:30}
\end{figure}

هناك صورة مكافئة متعددة الجسيمات. فالفراغ متكاثف، ويغيّر اهتزاز شعاعي كثافة المتكاثف. ولذلك يشبه بوزون هيغز موجة انضغاط في ذلك الوسط الكمومي. وهذه المماثلة مفيدة شريطة ألا يُستنتج منها أن الموجة لا بد أن تكون مستقرة. فموجات الصوت العادية يمكن أن تضمحل، كما يمكن لإثارة هيغز أن تضمحل أيضًا.

وهذا التمييز يفسّر لماذا كان العثور على بوزون هيغز القطعة الأخيرة المفقودة، لا أول دليل على الآلية. فقد كان الفراغ المنكسر والبوزونات الضعيفة الضخمة ظاهرين بالفعل على نحو غير مباشر. وكان لا بد من إنتاج موجة الكثافة الشعاعية على هيئة جسيم والتعرّف إليها من خلال نواتج اضمحلالها.

\section{الاضمحلال والإنتاج وميزة الفرميونات الثقيلة}

يقترن بوزون هيغز بفرميون اقترانًا يتناسب مع كتلة ذلك الفرميون. وقد يضمحل إلى زوج فرميون-ضد فرميون،
\begin{equation}
H\longrightarrow f+\bar f,
\end{equation}
ويمكن قراءة المخطط نفسه بالعكس بوصفه إنتاجًا،
\begin{equation}
f+\bar f\longrightarrow H.
\end{equation}
ولذلك فإن فناء إلكترون-بوزيترون ممكن من حيث المبدأ. لكن تاريخيًا افتقرت المصادمات المبكرة أولًا إلى طاقة كافية لإنتاج بوزون هيغز، ثم واجهت الآلات اللاحقة المشكلة الأعمق المتمثلة في أن اقتران يوكاوا للإلكترون صغير للغاية. وفناء الكواركات الخفيفة ضعيف بالمثل.

\begin{figure}[tbp]
\centering
\includegraphics[width=0.95\textwidth]{lecture_144_figure_12.png}
\caption{رأس اضمحلال هيغز المرسوم على اللوح. وتعطي قراءة الأسهم بالعكس اندماج فرميون وضد فرميون لتكوين بوزون هيغز.}
\lectureframe{00:59:40}
\end{figure}

كلما كان الفيرميون أثقل، كان اقترانه بيوكاوا أكبر. وهذا يميز كوارك القمة، الذي تبلغ كتلته نحو $173\,\mathrm{GeV}$، أي ما يقارب $180$ كتلة بروتون. وتستخدم المحاضرة الرقم التقريبي الدائري وهو نحو $170$ كتلة بروتون. إلا أن بوزون هيغز ذي كتلة تقارب $125\,\mathrm{GeV}$ أخف من أن يضمحل إلى زوج $t\bar t$ على الغلاف الكتلي،
\begin{equation}
m_H<2m_t.
\end{equation}
ولا يمكن أيضًا تكوين حزمة نافعة من كواركات القمة: إذ يضمحل كوارك القمة قبل أن يتمكن من التهادرن إلى جسيم حزمي طويل العمر.

توفر ميكانيكا الكم طريقًا لتجاوز كلتا العقبتين. يمكن لغلوونين أن يتقلبا عبر حلقة افتراضية لكوارك القمة وأن يخرجا على هيئة بوزون هيغز،
\begin{equation}
g+g\longrightarrow H.
\end{equation}
ولا يلزم لخطوط القمة الافتراضية أن تستوفي عتبة الغلاف الكتلي التي تحكم إنتاج أزواج القمة الحقيقية. وتمتلئ البروتونات بالغلوونات، لذا فإن تصادمات البروتون-البروتون عالية الطاقة في مصادم الهادرونات الكبير توفر مصدرًا كثيفًا لأزواج الغلوونات المطلوبة.

\begin{figure}[tbp]
\centering
\includegraphics[width=0.95\textwidth]{lecture_144_figure_13.png}
\caption{اندماج الغلوونات عبر حلقة افتراضية لكوارك القمة. يلتقي غلوونان مأخوذان من البروتونات المتصادمة في الحلقة ويُنتجان بوزون هيغز.}
\lectureframe{01:06:35}
\end{figure}

ولهذا السبب أصبح مصادم البروتونات، على الرغم من تعقيد حطام البروتونات، المصنع الحاسم لبوزون هيغز. فالبروتون لا يقترن ببوزون هيغز بوصفه جسمًا أوليًا واحدًا؛ بل إن مكوناته الغلوونية تُنتج بوزون هيغز عبر حلقة يهيمن عليها الاقتران بكوارك القمة الثقيل.

\section{جسيم عام 2012 واختباره بثنائي الفوتون}

كانت كتلة الجسيم المُعلَن عنه في 2012 قريبة من
\begin{equation}
m_H\simeq125\,\mathrm{GeV}.
\end{equation}
وفي المقارنة الميسرة في المحاضرة، يعادل ذلك تقريبًا $127$ كتلة بروتون. وقد أكمل اكتشافه قائمة الجسيمات التي يتطلبها النموذج المعياري الأدنى. وكان السؤال المتبقي هو ما إذا كانت خصائصه المقدارية مطابقة تمامًا لخواص النموذج المعياري.

يمكن قلب منطق حلقة القمة. يقترن بوزون هيغز بالجسيمات المشحونة الافتراضية في حلقة، ويمكن لتلك الحلقة أن تُصدر فوتونين،
\begin{equation}
H\longrightarrow\gamma+\gamma.
\end{equation}
هذا الاضمحلال نادر، لكن الفوتونين جسمان نظيفان على نحو استثنائي في مصادم هادروني. لذلك فإن زيادة ضيقة في توزيع الكتلة اللامتغيرة لهما يمكن أن تكشف عن بوزون هيغز حتى عندما تكون الاضمحلالات الأكثر شيوعًا أشد فوضى تجريبيًا.

\begin{figure}[tbp]
\centering
\includegraphics[width=0.95\textwidth]{lecture_144_figure_14.png}
\caption{مثلث ثنائي الفوتون على يمين اللوح، مقترنًا بحلقة اندماج الغلوونات السابقة على اليسار. تستقصي العمليتان المنطق نفسه على مستوى الحلقات ولكن في اتجاهين متعاكسين.}
\lectureframe{01:09:00}
\end{figure}

في وقت المحاضرة، بدا المعدل المُبلَّغ عنه لثنائي الفوتون معادلًا ربما مرةً ونصفَ التوقعَ الأبسطَ للنموذج المعياري، مع دلالة إحصائية لا تتجاوز نحو انحرافين معياريين. ولم يكن مثل هذا التذبذب دليلًا بحد ذاته. ولكن لو استمر، لكان يمكن أن يشير إلى جسيم مشحون إضافي في الحلقة، ربما كان مرتبطًا بالتناظر الفائق أو بامتداد آخر للنموذج المعياري.

تلك العبارة تنتمي إلى اللحظة التاريخية لـ 2012. وهي تسجل ما كان يُطلب من التجريبيين ترقبه بعد ذلك، لا ادعاءً راهنًا بوجود شذوذ. وكان الاكتشاف نفسه، إلى حد كبير، تأكيدًا لإطار كانت تدعمه بالفعل بيانات البوزونات الضعيفة. أما الاختبارات الأدق علميًا فكانت معدلات الاضمحلال: هل يقترن الجسيم السلمي الجديد بكل جسيم وفق النمط المعتمد على الكتلة الذي تتنبأ به آلية هيغز؟

\section{أسئلة من الحضور}

\begin{classroomqa}
\audiencequestion{ما الذي يحدد لماذا يتذبذب أحد الفرميونات بين اليسار واليمين أسرع من آخر، وبالتالي تكون له كتلة مختلفة؟\lecturetimestamp{01:12:00}}
\lecturerresponse{لكل نوع من الفرميونات اقترانه الخاص بمتكاثف هيغز. ويؤدي الاقتران الأقوى إلى مزج أقوى بين اليسار--اليمين وإلى كتلة أكبر. ويحدد الاقتران نفسه مدى قوة اضمحلال هيغز الفيزيائي إلى ذلك الفرميون، ولهذا تقترن الفرميونات الأثقل عمومًا على نحو أقوى. ويُدخل النموذج القياسي اقترانات يوكاوا هذه على هيئة معلمات، لكنه لا يفسر تدرجها الهرمي اللافت. والقول إن للإلكترون اقترانًا صغيرًا وإن لكوارك القمة اقترانًا كبيرًا هو وصف دقيق، لا اشتقاق أعمق بعد.}
\end{classroomqa}

\begin{classroomqa}
\audiencequestion{هل نعرف مقدار مجال هيغز في الفراغ؟\lecturetimestamp{01:13:36}}
\lecturerresponse{نعم. في الاصطلاح المستخدم هنا، قيمة التوقع في الفراغ هي نحو $240\,\mathrm{GeV}$؛ والقيمة المعيارية الأدق تقارب $246\,\mathrm{GeV}$. ويمكن تصورها إما على أنها الإزاحة الشعاعية للمجال عن الأصل، أو على أنها كثافة المتكاثف. وكمّ هيغز الفيزيائي هو تذبذب شعاعي صغير لتلك القيمة، أو على نحو مكافئ تقلب موضعي في كثافة المتكاثف.}
\end{classroomqa}

\begin{classroomqa}
\audiencequestion{هل تستلزم قيمة مجال هيغز غير المنعدمة في الفراغ كثافة طاقة للفضاء الخالي؟\lecturetimestamp{01:14:12}}
\lecturerresponse{نعم. يسهم كمون هيغز في طاقة الفراغ، كما تفعل قطاعات كثيرة أخرى من نظرية الحقل الكمي. وصافي طاقة الفراغ المرصود صغير على نحو استثنائي مقارنة بالمقادير الطبيعية لتلك الإسهامات، لذا فلا بد أن تلغي بعضها بعضًا تقريبًا. ولا توجد نظرية مقبولة تفسر هذا الإلغاء. وهذه هي مشكلة الثابت الكوني، وهي إحدى أعمق المشكلات غير المحلولة في الفيزياء الأساسية.}
\end{classroomqa}

\section{ما الذي حققه هذا البناء}

يمكن الآن قراءة هذا المنطق من البداية إلى النهاية من غير استدعاء أثير لزج. ويوفر الدوران الداخلي المُكمّى نموذجًا للشحنة. ويعطي تحويل الوعاء إلى قبعة مكسيكية للمجال قيمة غير منعدمة في الفراغ. ويحوِّل مبدأ عدم اليقين فراغًا كلاسيكيًا مُختارًا على نحو حاد إلى متكاثف غير حساس للكمّات المفردة. ويمكن للمتكاثف عندئذٍ أن يحمل فرق العدد الكمي الضعيف بين الفرميونات اليسارية واليمنية، مما يتيح امتزاجها وبالتالي أن تكون لها كتل. ويوفر أيضًا الأنماط الطولية التي تسمح لكل من $W$ و$Z$ بأن ينتشر بوصفه بوزونًا متجهًا ذا كتلة، بينما يبقى الفوتون عديم الكتلة.

بوزون هيغز الفيزيائي هو التذبذب الشعاعي الذي يبقى بعد إعادة الترتيب هذه. إنه ليس ``الكتلة نفسها،'' وليس مسؤولًا عن معظم كتلة البروتون. إنه موجة الكثافة القابلة للرصد لحقل كسر التناظر. وتحتفظ اقتراناته بأثر الكتل التي ولدها ذلك الحقل، ولهذا تهيمن الجسيمات الثقيلة على إنتاجه واضمحلاله. ولذلك فحلقة اندماج الغلوونات، وقناة الفوتونين، وسؤال الحضور الأخير عن طاقة الفراغ، ليست موضوعات منفصلة مثيرة للفضول. بل هي نتائج واختبارات وتبعات غير مكتملة للفراغ غير التافه نفسه.
